\section{Research Hypotheses}
\label{sec:research_objectives_research_hypothesis}

The three hypotheses correspond directly to the three research goals above.

\begin{itemize}

    \item \textbf{H1 (The Composability Gap Hypothesis).}
    We hypothesise that when pretrained generative models are logically composed via score addition, semantic composition is most reliable when the constituent concepts can be realised jointly without distorting one another. 
    When concepts overlap, compete for the same structure, or impose joint constraints on the output, logical composition diverges from semantic composition.
    These failures are not arbitrary: each class of failure exhibits predictable and characterisable failure modes, such as hybridisation, broken relations, or structural incoherence. 
    This results in a measurable divergence in latent trajectories, where logical composition enters regions of the manifold that are semantically unreachable under a single integrated prompt

    \item \textbf{H2 (The Representation Hypothesis).}
    We hypothesise that the composability gap is primarily a representation problem
    rather than an inference problem.
    Specifically, entangled latent representations cause separately composed concepts
    to interfere, and this cannot be corrected fully by modifying the inference-time
    composition rule alone.
    If this is correct, then models trained with explicit independence constraints
    should reduce or close the gap, whereas models without such constraints should
    continue to exhibit it even when composition rules are improved.

    \item \textbf{H3 (The Clinical Independence Sufficiency Hypothesis).}
    We hypothesise that approximate clinical independence between disease conditions
    is sufficient for factorised composition to produce ecologically valid outputs.
    If the disease-specific variation in a medical image can be approximately separated
    from shared anatomical and acquisition variation, then composing per-pathology
    models will produce outputs that correctly represent each condition, including
    co-morbid combinations not present in the training data.

\end{itemize}
